\documentclass[english, parskip=full]{scrartcl}
\usepackage[utf8]{inputenc}  % Kodierung der Datei
\usepackage[english]{babel}  % Sprache auf Deutsch 
\usepackage{color}
\usepackage{graphicx, gensymb}
\usepackage{amsfonts, cancel, amssymb, slashed}
\usepackage{amsmath, amsthm, array, esvect, esint}
\usepackage{booktabs, multirow}  % schönere Tabellen
\usepackage{caption}  % Anpassung der Captions
\usepackage{scrlayer-scrpage}    % Manipulation des Seitenstils
\pagestyle{scrheadings}  % Stil für die Seite setzen
\automark{section}  % Abschnittsnamen als Seitenbeschriftung 
\setheadsepline{.5pt}    % Kopzeile durch Linie abtrennen
\usepackage[hidelinks]{hyperref}  % Links und weitere PDF-Features
\usepackage{typearea, tikz, pgffor, verbatim}
\usetikzlibrary{calc, quotes}
\usepackage[cache=false, outputdir=build]{minted}
\usepackage{dsfont}


\definecolor{bg}{rgb}{0.9,0.9,0.9}
\newcommand\numberthis{\addtocounter{equation}{1}\tag{\theequation}}
\newcommand{\bra}{}
\newcommand{\kom}{}
\newcommand{\brak}{}
\newcommand{\braket}{}
\newcommand{\intx}{}
\newcommand{\intr}{}
\newcommand{\kla}{}
\newcommand{\klam}{}
\newcommand{\klammer}{}
\newcommand{\oper}{}
\newcommand{\tecxt}{}
\newcommand{\Bra}{}
\newcommand{\Ket}{}
\def\I{\text{i}}
\def\e{\text{e}}

\newcommand*{\titel}{09 Canonical quantization}
\newcommand*{\autor}{Joachim Schwardt}

\hypersetup{pdfauthor={\autor}, pdftitle={\titel}}

\subject{Quantum theory 2}
\title{\titel}
\author{\autor}
\date{\today}

\begin{document}

%\begin{titlepage}
\maketitle  % Titel setzen
%\tableofcontents  % Inhaltsverzeichnis setzen
%\end{titlepage}

\section{Questions}
The relation between $\Psi(x)$ and $a_p$ is
\begin{align}
\Psi(x)&=\mathcal{F}[a_p](x)=(2\pi)^{-3/2}\intr\int_{\mathbb{R}^{3}}{a_p\e^{\I px}}\,\mathrm{d}^{3}p.
\end{align}
The Lagrangian is $\mathcal{L}=\I|\psi|^2-\frac{1}{2m}|\nabla\psi|^2$.
The Lagrangian and Hamiltonian for the hydrogen atom are
\begin{align}
L&=\frac{m}{2}\dot{q}^2+\frac{\alpha}{|q|},\\
H&=\frac{p^2}{2m}-\frac{\alpha}{|q|},
\end{align}
where $p=\partial_{\dot{q}}L=m\dot{q}$.
The Lagrangian and Hamiltonian for the one-dimensional harmonic oscillator are
\begin{align}
L&=\frac{m}{2}\dot{q}^2-\frac{m\omega^2}{2}q^2,\\
H&=\frac{p^2}{2m}+\frac{m\omega^2}{2}q^2.
\end{align}
Key relationship?\newpage
\section{From Lagrangian to Hamiltonian}
Table of the given Lagrangian, the momentum $p:=\partial_{\dot{q}} L$, the inverse relationship $\dot{q}=\dot{q}(p)$ and the resulting Hamiltonian $H=p\dot{q}-L$.
\begin{table}[h!]
\centering
\begin{tabular}{c|c|c|c}
$L$&$p$&$\dot{q}(p)$&$H$\\ \hline
$\frac{m}{2}\dot{q}^2-\frac{\kappa^2}{2}q^2$&$m\dot{q} $&$\frac{p}{m}$&$\frac{p^2}{2m}+\frac{\kappa^2}{2}q^2$\\ \hline
$\frac{\alpha}{2}\dot{q}^2q^2$&$\alpha q^2\dot{q} $&$\frac{p}{\alpha q^2}$&$\frac{p^2}{2\alpha q^2}$\\ \hline
$\left( \alpha\dot{q}+\beta q \right)^2$&$2\alpha\left( \alpha\dot{q}+\beta q \right)$&$\frac{1}{\alpha}\left( \frac{p}{2\alpha}-\beta q \right)$&$\frac{p^2}{4\alpha^2}-\frac{\beta}{\alpha}pq$\\ \hline
$\frac{m}{2}\dot{q}^2+Q\dot{q}\cdot A-Q\phi$&$m\dot{q} +QA$&$\frac{p-QA}{m}$&$\frac{1}{2m}(p-QA)^2+Q\phi$\\ 
\end{tabular}
\end{table}\\
Note that the usual notation for vectors was omitted in the last example.\newpage
\section{Linear chain}
Let the coordinates $q_n$ define a periodic lattice with lattice distance $a$, i.e. $q_{n+N}\equiv q_n$ for all $n=1,\dots,N$.
The total length is $L:=Na$.
The Lagrangian is
\begin{align}
L&=\sum_{n=1}^N\frac{\dot{q}_n^2}{2}-\frac{\kappa}{2}(q_{n+1}-q_n)^2,
\end{align}
where the mass is assumed to be $m=1$.
The Euler-Lagrange equations are 
\begin{align}
0&=\frac{\tecxt\text{d}}{\tecxt\text{d}t}\partial_{\dot{q}_j}L-\partial_{q_j}L=\sum_{n=1}^N\ddot{q}_n\delta_{jn}+\kappa (q_{n+1}-q_n)(\delta_{j,n+1}-\delta_{jn})\\
\Rightarrow\ \ \ \ddot{q}_j&=-\kappa (q_j-q_{j-1})+\kappa (q_{j+1}-q_j)=\kappa (q_{j+1}+q_{j-1}-2q_j).
\end{align}
Let $q_n(t)=\e^{\pm\I (\omega_kt-kna)}$.
Then the periodic boundary condition yields
\begin{align}
q_{n+N}&=\e^{\mp\I kNa}q_n\overset{!}{=}q_n,\\
\Rightarrow\ \ \ \e^{\mp\I kNa}&=1.
\end{align}
This is equivalent to $kNa=2\pi z$ for $z\in\mathbb{Z}$.
The minimal non-zero value is $|k|=\frac{2\pi}{Na}=\frac{2\pi}{L}$.
Due to the periodic condition it is sufficient to choose $-\frac{N}{2}<z\le\frac{N}{2}$, or equivalently
\begin{align}
-\pi &< ka\le \pi.
\end{align}
This range contains exactly $N$ values. 
Since $N$ is assumed to be odd, we can write $N=2M+1$ with $M\in\mathbb{N}$.
Hence the possible $z$ are $z=-M,-M+1,\dots M$.
The largest wavenumber is 
\begin{align}
k_\tecxt\text{max}&=\frac{2\pi}{L}z_\tecxt\text{max}=\frac{2\pi M}{Na}=\frac{\pi}{a}\frac{2M}{2M+1}=\frac{\pi}{a}-\frac{\pi}{L}.
\end{align}
Inserting the Ansatz
\begin{align}
-\frac{\omega_k^2}{\kappa}q_n&=\kla\left( {\e^{\mp\I ka}+\e^{\pm\I ka}-2} \right)q_n=2\kla\left( {\cos (ka)-1} \right)q_n
\end{align} 
leads to the dispersion relation
\begin{align}
\omega_k&=2\sqrt{\kappa}\left\vert\sin\frac{ka}{2}\right\vert=2\sqrt{\kappa}\left\vert\sin\frac{\pi z}{N}\right\vert.
\end{align}
In the limit $k\rightarrow 0$ we can approximate this relation by
\begin{align}
\omega_k&\simeq \sqrt{\kappa}ka.
\end{align}
For fixed $L$ and $N\rightarrow\infty$ with $a\rightarrow 0$ we find a continuous variable $x:=\frac{z}{N}\in \klam\left[ {-\frac{1}{2},\frac{1}{2}} \right]$.
The dispersion relation becomes
\begin{align*}
\omega (x) &:=2\sqrt{\kappa}\left\vert\sin (\pi x)\right\vert\le 2\sqrt{\kappa}.
\end{align*}
FIXME: this is just a guess, more or less unmotivated...\\
The speed of sound is given by 
\begin{align}
v&=Lf_\tecxt\text{max}=L\frac{\omega_\tecxt\text{max}}{2\pi}=\frac{L\sqrt{\kappa}}{\pi}.
\end{align}
The canonical momenta are $p_n=\dot{q}_n$.
The Hamiltonian then reads
\begin{align}
H&=\sum_{n=1}^N\frac{p^2}{2}+\frac{\kappa}{2}(q_{n+1}-q_n)^2.
\end{align}
\subsection{Canonical quantization}
Let $k=k(z)=\frac{2\pi}{Na}z$ and write $b_z:=a_z \e^{-\I (\omega (z)t -k(z)an)}$.
We set $b_z^\dagger := (b_z)^\dagger$ and define the Ansatz
\begin{align}
q_n&=\sum_{z=-M}^M\frac{1}{\sqrt{2N\omega}}\kla\left( {b_z+b_z^\dagger} \right).
\end{align}
We find the momentum operators
\begin{align}
p_n&=\dot{q}_n=\sum_{z=-M}^M\frac{-\I\omega}{\sqrt{2N\omega}}\kla\left( {b_z-b_z^\dagger} \right).
\end{align}
Consider the commutation relation
\begin{align}
\klam\left[ {q_n,q_m} \right]&=\frac{1}{2N}\sum_{z,x}\frac{1}{\sqrt{\omega_z\omega_x}}\klam\left[ {b_z^{(n)}+b_z^{(n)\dagger},b_x^{(m)}+b_x^{(m)\dagger}} \right]\\
&=\frac{1}{2N}\sum_{z,x}\kla\left( {\delta_{zx}-\delta_{zx}} \right)=0.
\end{align}
Here we assume the usual relations $\klam\left[ {a_z,a_x^\dagger} \right]=\delta_{zx}$ and $\klam\left[ {a_z,a_x} \right]=0$, which automatically lead to
\begin{align}
\klam\left[ {b_z^{(n)},b_x^{(m)\dagger}} \right]&=\e^{-\I (\omega (z)t -k(z)an)}\e^{\I (\omega (x)t -k(x)am)}\klam\left[ {a_z,a_x^\dagger} \right]=\e^{\I k(z)a(n-m)}\delta_{zx},\\
\klam\left[ {b_z^{(n)},b_x^{(m)}} \right]&=\e^{-\I (\omega (z)t -k(z)an)}\e^{-\I (\omega (x)t -k(x)am)}\klam\left[ {a_z,a_x} \right]=0.
\end{align}
This allows us to compute the commutator
\begin{align}
\klam\left[ {q_n,p_m} \right]&=\frac{-\I}{2N}\sum_{z,x}\sqrt{\frac{\omega (x)}{\omega (z)}}\klam\left[ {b_z^{(n)}+b_z^{(n)\dagger},b_x^{(m)}-b_x^{(m)\dagger}} \right]=\\
&=\frac{\I}{2N}\sum_{z,x}\sqrt{\frac{\omega (x)}{\omega (z)}}\kla\left( {\e^{\I k(z)a(n-m)}+\e^{\I k(z)a(m-n)}} \right)\delta_{zx}=\\
&=\frac{\I}{N}\sum_{z=-M}^M\e^{\I 2\pi\frac{z}{N}(n-m)}=\I\delta_{nm}.
\end{align}
In the last step we change $z\rightarrow -z$ for the summation over the second exponential phase, which eliminates the factor $\frac{1}{2}$, as both exponential factors become identical.
We also use the relation
\begin{align}
S_{nm}&:=\sum_{z=-M}^M\e^{\I 2\pi\frac{z}{N}l}\kom\overset{\substack{{x\,:=\,z+M} \\ |}}{=}\\
&=\sum_{x=0}^{N-1}\kla\left( {\e^{\I 2\pi \frac{l}{N}}} \right)^x=\\
&=\frac{1-\e^{\I 2\pi l}}{1-\e^{\I 2\pi \frac{l}{N}}}=N\delta_{l,0}.
\end{align}
Here $l\in\mathbb{Z}$ is any integer, and we use $N=2M+1$.
Above we had $l=n-m$.
\\
The Hamiltonian is
\begin{align}
H&=\sum_{n=1}^N\frac{p_n^2}{2}+\frac{\kappa}{2}(q_{n+1}-q_n)^2=:T+V.
\end{align}
The kinetic term is given by
\begin{align}
T&=\frac{1}{2}\sum_{n=1}^Np_n^2=\frac{(-\I)^2}{4N}\sum_{n,z,x}\sqrt{\omega (z)\omega (x)}\kla\left( {b_zb_x+b_z^\dagger b_x^\dagger -b_zb_x^\dagger -b_z^\dagger b_x} \right)=\\
&=\frac{1}{4N}\sum_{n,z}\omega (z)+\frac{1}{4N}\sum_{n,z,x}\sqrt{\omega (z)\omega (x)}\kla\left( {b_z^\dagger b_x+b_x^\dagger b_z-b_z b_x-b_z^\dagger b_x^\dagger} \right)=\\
&=\frac{\sqrt{\kappa}}{2}\sum_z\left\vert\sin\frac{\pi z}{N}\right\vert+\frac{1}{4N}\sum_{z,x}\sqrt{\omega (z)\omega (x)}a_z^\dagger a_x\sum_{n=1}^N\e^{\I 2\pi\frac{n}{N}(x-z)}+\dots=\\
&=\tecxt\text{const}+\frac{1}{2}\sum_{z=-M}^M\omega_za_z^\dagger a_z-\frac{1}{4N}\sum_{z,x}\sqrt{\omega_z\omega_x}a_za_x\sum_{n=1}^N\e^{-\I \kla\left( {[\omega_z+\omega_x]t-[k_z+k_x]an} \right)}+\dots=\\
&=\tecxt\text{const}+\frac{1}{2}\sum_k\omega_k\hat{n}_k-\frac{1}{4N}\sum_{z,x}\sqrt{\omega_z\omega_x}a_za_x\e^{-\I[\omega_z+\omega_x]t}\sum_{n=1}^N\e^{\I 2\pi \frac{n}{N}(z+x)}+\dots=\\
&=\tecxt\text{const}+\frac{1}{2}\sum_k\omega_k\hat{n}_k-\frac{1}{4}\sum_z\omega_z\kla\left( {a_za_{-z}\e^{-2\I\omega_zt}+a_z^{\dagger}a_{-z}^\dagger\e^{2\I\omega_zt}} \right).
\end{align}
The potential term is more complicated, but leads to a very similar result.
First consider the difference
\begin{align}
\sqrt{\kappa}(q_{n+1}-q_n)&=\sum_{z=-M}^M\frac{\sqrt{\kappa}}{\sqrt{2N\omega (z)}}\klam\left[ {b_z\kla\left( {\e^{\I k(z)a}-1} \right)+b_z^\dagger\kla\left( {\e^{-\I k(z)a}-1} \right)} \right]=\\
&=\sum_{z=-M}^M\frac{\sqrt{\kappa}}{\sqrt{2N\omega (z)}}\klam\left[ {b_z\e^{\I k(z)a/2}2\I\sin\frac{k(z)a}{2}\pm b_z^\dagger \e^{-\I k(z)a/2}2\I\sin\frac{\mp k(z)a}{2}} \right]=\\
&=\sum_{z=-M}^M\frac{\I\omega}{\sqrt{2N\omega}}\klam\left[ {b_z\e^{\I k(z)a/2}\pm b_z^\dagger\e^{-\I k(z)a/2}} \right].
\end{align}
Notice the ambiguity of the sign for the second term. 
This is due to $\omega_k\sim |\sin(ka)|$.
A more careful derivation shows, i.e. without inserting the dispersion relation before squaring the term, that the $+$-sign is correct here.
\begin{align}
V&=\frac{\I^2}{4N}\sum_{n,z,x}\sqrt{\omega_z\omega_x}\kla\left( {b_zb_x\e^{\I [k_z+k_x]a/2}+\dagger +b_zb_x^\dagger\e^{\I \Delta k_{zx} a/2} +b_z^\dagger b_x\e^{-\I \Delta k_{zx} a/2}} \right)=\\
&=\frac{1}{4N}\sum_{n,z}\omega_z+\frac{1}{4N}\sum_{n,z,x}\sqrt{\omega_z\omega_x}\kla\left( {b_z^\dagger b_x\e^{-\I \Delta k_{zx} a/2}+\dagger+b_z b_x\e^{\I [k_z+k_x]a/2}+\dagger} \right)=\\
&=\tecxt\text{const}+\frac{1}{4N}\sum_{z,x}\sqrt{\omega_z\omega_x}a_z^\dagger a_x\e^{-\I \Delta k_{zx} a/2}\e^{-\I\Delta\omega_{zx}t}\sum_{n=1}^N\e^{\I 2\pi\frac{n}{N}(x-z)}+\dots=\\
&=\tecxt\text{const}+\frac{1}{2}\sum_{z=-M}^M\omega_za_z^\dagger a_z+\frac{1}{4N}\sum_{z,x}\sqrt{\omega_z\omega_x}a_za_x\e^{\I [k_z+k_x]a/2}\e^{-\I[\omega_z+\omega_x]t}\sum_{n=1}^N\e^{-\I[k_z+k_x]an}+\dagger=\\
&=\tecxt\text{const}+\frac{1}{2}\sum_k\omega_k\hat{n}_k+\frac{1}{4}\sum_z\omega_z\kla\left( {a_za_{-z}\e^{-2\I\omega_zt}+a_z^{\dagger}a_{-z}^\dagger\e^{2\I\omega_zt}} \right).
\end{align}
Combining both expressions yields
\begin{align*}
H&=\sum_k\omega_k\hat{n}_k+\frac{1}{2}\sum_k\omega_k.
\end{align*}
%const=\sqrt{\kappa}\sum_{z=-M}^M\e^{\I\pi \frac{z}{N}}=\frac{2}{1-\e^{\I\pi/N}}
\end{document}